Damit Schweiss auf der Haut verdunsten kann, muss ihm die Verdampfungswärme zugeführt werden; bei Körpertemperatur beträgt sie etwa \LvWO\ (FoTaBe S.~193). Diese Wärme wird zum grössten Teil der Haut und damit dem Körper entzogen. Jedes Gramm verdunsteter Schweiss kühlt den Körper also beträchtlich, ohne dass die Umgebung kälter sein muss als die Haut. So wird die erhöhte Körpertemperatur wieder gesenkt.
